\documentclass[10pt,a4paper]{amsart}
\usepackage[margin=19mm]{geometry}
\usepackage{amsmath,amssymb,mathtools}
\usepackage[scr=rsfs,cal=euler]{mathalfa}
\usepackage{xfrac}
\usepackage[unicode,hidelinks,bookmarks=false]{hyperref}
\newcommand{\ie}{i.e.\ }
\newcommand{\cf}{cf.\ }
\newcommand{\resp}{respectively\ }
\newcommand{\Subsection}{\subsection}
\providecommand{\nobreakdash}{\nobreak\hbox{-}}
\setlength{\emergencystretch}{2em}
\multlinegap0pt
\usepackage{bm}
\usepackage{bbm}
\usepackage{dsfont}
\usepackage{xspace}
\usepackage{cancel}
\usepackage{mathtools}
\usepackage{amsthm}
\makeatletter
\@ifundefined{theorem}{\newtheorem{theorem}{Theorem}}{}
\@ifundefined{definition}{\newtheorem{definition}[theorem]{Definition}}{}
\@ifundefined{lemma}{\newtheorem{lemma}[theorem]{Lemma}}{}
\makeatother
\newcommand{\Div}[2]{#1\,\mid\, #2}
\newcommand{\Mod}[3]{#1\equiv #2 \pmod{#3}}
\newcommand{\Nat}{\mathbb{N}}
\newcommand{\Set}[2]{\set{#1\, :\, #2}}
\newcommand{\nDiv}[2]{#1\,\nmid\, #2}
\newcommand{\set}[1]{\left\{#1\right\}}
\title[Counting degrees of vertices in near Goldbach graphs]{Counting degrees of vertices in near Goldbach graphs}
\author{Shamik Ghosh, Souradeep De}
\date{Source: arXiv:2608.14159v2 (CC BY 4.0)}
\begin{document}
\maketitle

\noindent\emph{Source and licence.} Counting degrees of vertices in near Goldbach graphs. Version arXiv:2608.14159v2 (CC BY 4.0), distributed under the Creative Commons Attribution 4.0 International licence, as recorded in the arXiv licence metadata for this exact version and shown on the abstract page https://arxiv.org/abs/2608.14159v2. Full rights notice: licenses/arxiv-2608.14159-CC-BY-4.0.txt.\par\smallskip

\noindent\textit{Editorial scope.} Counted unit: the Lemma whose source label is \texttt{lem1good} in the article (source0.tex lines 1044--1046, source label \texttt{lem1good}), delivered together with the article's own complete original proof of it (source0.tex lines 1048--1087, one \texttt{proof} environment of 40 lines). No mathematics is written, repaired or re-derived: statement and proof are the article's own text line for line. Required context retained uncounted: eqfn (lines 334--336); lemfp4 (lines 350--358); defap (lines 390--393); defeta13 (lines 399--422); lemins (lines 1009--1011); lemk (lines 1021--1024). Every definition, display and label of those lines is retained unchanged. Omitted from this package and not redistributed: 1--333, 337--349, 359--389, 394--398, 423--1008, 1012--1020, 1025--1043, 1047--1047, 1088--1546. Nothing inside a retained excerpt is omitted or altered: every retained body is delivered exactly as the article's lines are written, so no omission receipt of any kind accompanies this package. Citations: the retained excerpts carry no citation command at all, so no bibliography entry of the article is transcribed and no reference is redistributed. Reference adaptation: none was needed and none was made; every reference inside the delivered unit resolves inside it, so no unresolved reference or citation is produced. Printed slips kept literal: no printed word, symbol, hypothesis or number of the counted statement or of its proof was altered. Layout and class: the article's own class is \texttt{article} with packages including babel, inputenc, todonotes, float, amsmath, amssymb, amsthm, cleveref, graphicx, mathrsfs, multicol, footmisc; the wrapper is the lane's pinned \texttt{amsart} editorial wrapper at 10pt on A4, which restates the theorem-like environments the retained text needs and adds the article's own macro definitions that the retained text needs (\texttt{\textbackslash Div}, \texttt{\textbackslash Mod}, \texttt{\textbackslash Nat}, \texttt{\textbackslash Set}, \texttt{\textbackslash nDiv}, \texttt{\textbackslash set}), with extra wrapper packages (bm, bbm, dsfont, xspace, cancel, mathtools). The delivered numbering is the unit's own. Assets: no retained span contains a figure environment, a graphics-inclusion command, a PDF-page inclusion, a table or a TikZ picture; no image file is copied into this package, the package declares no asset dependency and no image file is redistributed with it. Licence: version arXiv:2608.14159v2 of this article is distributed under the Creative Commons Attribution 4.0 International licence, as recorded in the arXiv licence metadata for this exact version and shown on the abstract page https://arxiv.org/abs/2608.14159v2. A full rights notice is delivered at licenses/arxiv-2608.14159-CC-BY-4.0.txt. Characters: every retained excerpt is pure ASCII, so the delivered engine, XeLaTeX with its default Latin Modern text font, reports no missing character.
\begin{equation}\label{eqfn}
f(n)=|\bigcup\limits_{\substack{p\, =\, \text{prime}\\ 3<p<\sqrt{x}}} \ \left(F(x,p,+)\cup F(x,p,-)\right)|.
\end{equation}

\begin{lemma}\label{lemfp4}
$|F(x,p,+)|=\lfloor (n-k_1)/p\rfloor+1$, $|F(x,p,-)|=\lfloor (n-k_2)/p\rfloor $ if\ $\Div{3}{p+1}$ and $|F(x,p,-)|=\lfloor (n-k_2)/p\rfloor +1$ if\ $\nDiv{3}{p+1}$.  
$$|F(x,p,+)\cup F(x,p,-)|=\left\{\begin{array}{ll}
\lfloor (n-k_1)/p\rfloor+1, & \text{if } \Div{p}{x}\\
\lfloor (n-k_1)/p\rfloor + \lfloor (n-k_2)/p\rfloor +1, & \text{if } \nDiv{p}{x}\text{ and } \Div{3}{p+1}\\
\lfloor (n-k_1)/p\rfloor + \lfloor (n-k_2)/p\rfloor +2, & \text{if } \nDiv{p}{x}\text{ and } \nDiv{3}{p+1},\\
\end{array}\right.$$
where $k_1=[(p-[x]_p^1-6)\times 12^{p-2}]_p+1$, $k_2=[(x-6)\times 12^{p-2}]_p+1$, $[a]_p=i\in\set{0,1,2,\ldots,p-1}$ and $[a]_p^1=i\in\set{1,2,\ldots,p}$ such that $\Mod{a}{i}{p}$.
\end{lemma}

\begin{definition}\label{defap}
Let $n\in\Nat$, $x=12n+4$, $S(n)=\Set{12m+6}{m\in\Nat\cup\set{0},\ m<n}$.\\
 Let $PR(x)$ be the set of primes $p$ such that $3<p<\sqrt{x}$ and $r_x=|PR(x)|$. For each prime $p\in PR(x)$, let $A_{p}$ be the event which is the set of numbers $y$ in $S(n)$ such that either $\Div{p}{x+y}$ or $\Div{p}{x-y}$ (except the case where $p=\frac{x-y}{2}$). Let $\mathscr{A}(x)=\Set{A_p}{p\in PR(x)}$ and $\mathscr{A}_x=\bigcup \Set{A_p}{p\in PR(x)}$. We call $\mathscr{A}(x)$ as the {\em set of divisibility events} for $x$. Note that for $n>3$, $|\mathscr{A}(x)|=|PR(x)|>1$. Moreover, we define $T(n)=\frac{f(n)}{n}$, where $f(n)$ is given by (\ref{eqfn}).
\end{definition}

\begin{definition}\label{defeta13}
Let $n\in\Nat$ and $x=12n+4$. For any prime $p>3$, we define
$$pr_1(p)=\left\{\begin{array}{ll}
\frac{\lfloor (n-k_1)/p\rfloor+1}{n}, & \text{if } \Div{p}{x}\\
\frac{\lfloor (n-k_1)/p\rfloor + \lfloor (n-k_2)/p\rfloor +1}{n}, & \text{if } \nDiv{p}{x}\text{ and } \Div{3}{p+1}\\
\frac{\lfloor (n-k_1)/p\rfloor + \lfloor (n-k_2)/p\rfloor +2}{n}, & \text{if } \nDiv{p}{x}\text{ and } \nDiv{3}{p+1},\\
\end{array}\right.$$
where $k_1$ and $k_2$ are defined as in Lemma \ref{lemfp4} and

$$pr_2 (p)=\left\{\begin{array}{ll}
\frac{1}{p}+\frac{1}{n}, & \text{if } \Div{p}{x}\\
\frac{2}{p}+\frac{1}{n}, & \text{if } \nDiv{p}{x}\text{ and } \Div{3}{p+1}\\
\frac{2}{p}+\frac{2}{n}, & \text{if } \nDiv{p}{x}\text{ and } \nDiv{3}{p+1}.\\
\end{array}\right.$$
Now we define the following approximations of $T(n)$:

\begin{equation}\label{eqneta5n}
\tau_1(n)=1-\prod\limits_{i=1}^{r_x} (1-pr_1(p_i)).
\end{equation}

\begin{equation}\label{eqneta6n}
\tau_2(n)=1-\prod\limits_{i=1}^{r_x} (1-pr_2 (p_i)).
\end{equation}
\end{definition}

\begin{lemma}\label{lemins}
Let $A$ and $B$ be two independent events that occur with probabilities $P(A)$ and $P(B)$ respectively. Let $p,q\in [0,1]$ such that $P(A)<p$ and $P(B)\leq q$. Then $P(A\cup B)<1 - (1 - p) (1 -q)$.
\end{lemma}

\begin{lemma}\label{lemk}
Let $A$ and $B$ be two events that occur with probabilities $P(A)$ and $P(B)$ respectively\\[0.5em]
 such that $0<P(A),P(B)<1$ and $\frac{P(B)-P(A\cap B)}{P(B)-P(A)P(B)}<k$. Then $P(A\cup B)<1-(1-P(A))(1-kP(B))$.
\end{lemma}

\begin{lemma}\label{lem1good}
Let $n\in\Nat$, $n>17$, then $0<T(n)<\tau(n)<1$.
\end{lemma}

\begin{proof}
We first note that, since $1\leq k_1,k_2\leq p_i$, $pr_1(p_i)\leq \frac{2}{p_i}+\frac{2}{n}$ for all $p_i\in PR(x)$. Since $p_i\geq 7$ for all $i\geq 2$, we have $2.5\, pr_1(p_i)\leq \frac{5}{2}\, (\frac{2}{p_i}+\frac{2}{n})=\frac{5}{p_i}+\frac{5}{n}<1$ for $n>17$. So $(1-2.5\, pr_1(p_i))\in (0,1)$. Also $pr_1(5)\in (0,1)$. So $\tau(n)\in (0,1)$. 

\vspace{1em}
\noindent
Now $T(n)=\frac{f(n)}{n}$, where $f(n)=|\bigcup\limits_{\substack{p\, =\, \text{prime}\\ 3<p<\sqrt{x}}} \ \left(F(x,p,+)\cup F(x,p,-)\right)|$ by (\ref{eqfn}).

\vspace{1em}
\noindent
 Following Definition \ref{defap}, for any prime $p_i$, ($3<p_i<\sqrt{x}$),\\
 $P(A_{p_i})=|F(x,p_i,+)\cup F(x,p_i,-)|/n=pr_1(p_i)$ (see Lemma \ref{lemfp4} and Definition \ref{defeta13}) and 
$$T(n)=P(\bigcup\limits_{i=1}^{r_x} A_{p_i}).$$
We prove $T(n)<\tau(n)$ by induction. 

\vspace{1em}
\noindent
Let
$$P_{t}=1-(1-pr_1(5))\prod\limits_{i=2}^{t} (1-2.5\, pr_1(p_i)).$$
We have $P(A_5)=pr_1(5)=P_1$. Let $A=A_5$ and $B=A_7$. Now $\lambda_2=\lambda(A,B)\leq \lambda_0<2.5$. Then $P(A\cup B)<1-(1-P(A))(1-2.5\, P(B))$ by Lemma \ref{lemk} considering $k=2.5$.\\
 Thus $P(A_5\cup A_7)<1-(1-pr_1(5))(1-2.5\, pr_1(7))=P_2$. In fact, we can prove this directly as follows:
$$P(A_5\cup A_7)=P(A_5)+P(A_7)-P(A_5\cap A_7)\leq P(A_5)+P(A_7)=pr_1(5)+pr_1(7).$$

\noindent
Now $1-(1-pr_1(5))(1-2.5\, pr_1(7))=pr_1(5)+2.5\, pr_1(7)-2.5\, pr_1(5)\, pr_1(7)$\\[0.5em]
\null\hspace{2.3in} $= \left(pr_1(5)+pr_1(7)\right)+\left(1.5-2.5\, pr_1(5)\right)\, pr_1(7)$.\\[0.5em]
Also $pr_1(5)\leq \frac{2}{5}+\frac{2}{n}<0.6=\frac{1.5}{2.5}$ (for $n> 10$) $\Longrightarrow \left(1.5-2.5\, pr_1(5)\right)>0$\\[0.5em]
$\Longrightarrow 1-(1-pr_1(5))(1-2.5\, pr_1(7))>pr_1(5)+pr_1(7)\geq P(A_5\cup A_7) \text{ as }pr_1(7)>0.$
Thus we proved that 
$$P(\bigcup\limits_{i=1}^{2} A_{p_i}) <P_2.$$

\vspace{1em}
\noindent
Let $A=\bigcup\limits_{i=1}^{t-1} A_{p_i}$ and $B=P(A_{p_t})=pr_1(p_t)$ for some $t>2$. Suppose $P(A)<P_{t-1}$. Now $\lambda_t=\lambda(A,B)\leq \lambda_0<2.5$. Then 
$$P(A\cup B)<1-(1-P(A))(1-2.5\, P(B))$$
by Lemma \ref{lemk} considering $k=2.5$. Therefore 
$$P(\bigcup\limits_{i=1}^{t} A_{p_i})=P(A\cap B)<P(A)+2.5\, P(B)-2.5\, P(A)P(B)<1-(1-P_{t-1})(1-2.5\, pr_1(p_t))=P_t$$
by Lemma \ref{lemins} as $P(A)<P_{t-1}$ and $P(B)=pr_1(p_t)$. This completes the induction. Thus finally, we have 
$$T(n)=P(\bigcup\limits_{i=1}^{r_x} A_{p_i})<\tau(n).$$
Therefore, $0<T(n)<\tau(n)<1$ as required.
\end{proof}

\end{document}
